在下一个命题及其证明中，给定 $X$ 中从 $x$ 出发、或回到 $x$ 或走到 $\sigma(x)$ 的道路
$\gamma$（$\tilde\gamma$ 为其从 $\tilde x$ 出发的唯一提升），我们用
$\gamma\curvearrowright\tilde x$ 记元素 $\tilde\gamma(1)$。

\begin{proposition}\label{prop:2.2.17}
设 $x$ 是 $X$ 中满足 $\sigma(x)\ne x$ 的点。令 $\Gamma$ 为形如
\[
\{\zeta:[0,1]\to X\mid \zeta(0)=\zeta(1)=x\}
\sqcup\{\eta:[0,1]\to X\mid \eta(0)=x,\ \eta(1)=\sigma(x)\}
\]
的道路（端点固定）的同伦类之集，并对 $\Gamma$ 中所有 $\gamma,\gamma'$ 规定
\[
\gamma\boxtimes\gamma'=
\begin{cases}
\gamma*\gamma', & \text{若 } \gamma'(1)=x,\\[0.3em]
(\sigma\circ\gamma)*\gamma', & \text{若 } \gamma'(1)=\sigma(x).
\end{cases}
\]
由万有覆盖的泛性质与引理 2.2.16，给定 $\gamma\in\Gamma$，存在唯一的自微分同胚
$f_\gamma:\tilde X\to\tilde X$ 使得 $f_\gamma(\tilde x)=\gamma\curvearrowright\tilde x$。
若 $\gamma(1)=x$ 则 $f_\gamma$ 覆盖 $\Id_X$，若 $\gamma(1)=\sigma(x)$ 则它覆盖 $\sigma$。
于是 $(\Gamma,\boxtimes)$ 是群，且映射
\[
\Phi:\ \Gamma\longrightarrow\pi^R_1((X,\sigma);x),
\qquad \gamma\longmapsto f_\gamma
\]
是群同构；事实上它使下图可交换：
\[
\begin{tikzcd}[column sep=large, row sep=large]
1 \arrow[r] & \pi_1(X;x) \arrow[r] \arrow[d, "\Phi\vert_{\pi_1(X;x)}"', "\simeq"] & \Gamma \arrow[r, "\beta"] \arrow[d, "\Phi"', "\simeq"] & \Z/2\Z \arrow[r] \arrow[d, "="] & 1\\
1 \arrow[r] & \pi_1(X;x) \arrow[r] & \pi^R_1((X,\sigma);x) \arrow[r, "\alpha"'] & \Z/2\Z \arrow[r] & 1
\end{tikzcd}
\]
其中若 $\gamma(1)=x$ 则 $\beta(\gamma)=\Id_X$，若 $\gamma(1)=\sigma(x)$ 则
$\beta(\gamma)=\sigma(x)$。
\begin{annote}
底本此处作 $\beta(\gamma)=\sigma(x)$，与 $\beta$ 的取值落在
$\Z/2\Z=\{\Id_X,\sigma\}$ 不符；按前后文，第二式应为 $\beta(\gamma)=\sigma$。译文照排底本，
谨此注明。
\end{annote}
\end{proposition}

\begin{proof}
把 $\sigma\circ\gamma:[0,1]\to X$ 的道路（同伦类）简记为 $\sigma(\gamma)$ 会带来方便。
利用 $\sigma(\gamma*\gamma')=\sigma(\gamma)*\sigma(\gamma')$，不难验证 $\Gamma$ 上的乘积
$\boxtimes$ 满足结合律且所有元素都有逆（单位元是 $x$ 处常道路的同伦类）。故 $\Gamma$ 是
群。

设 $\gamma,\gamma'$ 是 $\Gamma$ 中两个元素，记 $\Phi(\gamma)=f_\gamma$、
$\Phi(\gamma')=f_{\gamma'}$。则 $\Phi(\gamma\boxtimes\gamma')$ 是唯一的映射
$h:\tilde X\to\tilde X$，它覆盖 $\Id_X$（若 $\gamma\boxtimes\gamma'$ 止于 $x$）或覆盖
$\sigma$（若 $\gamma\boxtimes\gamma'$ 止于 $\sigma(x)$），且满足
\[
h(\tilde x)=(\gamma\boxtimes\gamma')\curvearrowright\tilde x=
\begin{cases}
(\gamma*\gamma')\curvearrowright\tilde x, & \text{若 } \gamma'(1)=x,\\[0.3em]
(\sigma(\gamma)*\gamma')\curvearrowright\tilde x, & \text{若 } \gamma'(1)=\sigma(x).
\end{cases}
\]
令当 $\gamma'(1)=x$ 时 $\alpha:=\Id_X$，当 $\gamma'(1)=\sigma(x)$ 时 $\alpha:=\sigma$，
于是 $h(\tilde x)=(\alpha\circ\gamma)\curvearrowright(\gamma'\curvearrowright\tilde x)
=(\alpha\circ\gamma)\curvearrowright f_{\gamma'}(\tilde x)$。此时
$f_{\gamma'}\circ\tilde\gamma$ 是始于 $f_{\gamma'}(\tilde\gamma(0))
=f_{\gamma'}(\tilde x)$ 的道路，且满足
\[
p\circ(f_{\gamma'}\circ\tilde\gamma)=(p\circ f_{\gamma'})\circ\tilde\gamma
=(\alpha\circ p)\circ\tilde\gamma
=\alpha\circ(p\circ\tilde\gamma)=\alpha\circ\gamma.
\]
故 $(\alpha\circ\gamma)\curvearrowright f_{\gamma'}(\tilde x)
=(f_{\gamma'}\circ\tilde\gamma)(1)
=f_{\gamma'}(\tilde\gamma(1))
=f_{\gamma'}(\gamma\curvearrowright\tilde x)$，从而
\[
h(\tilde x)=(\alpha\circ\gamma)\curvearrowright f_{\gamma'}(\tilde x)
=f_{\gamma'}(\gamma\curvearrowright\tilde x)
=f_{\gamma'}\bigl(f_\gamma(\tilde x)\bigr)
=(f_{\gamma'}\circ f_\gamma)(\tilde x).
\]
此外，若 $h$ 覆盖 $\sigma$ 则 $f_{\gamma'}\circ f_\gamma$ 也覆盖 $\sigma$，若 $h$ 覆盖
$\Id_X$ 则它也覆盖 $\Id_X$。故由这种映射的唯一性
$\Phi(\gamma\boxtimes\gamma')=h=f_{\gamma'}\circ f_\gamma
=\Phi(\gamma')\circ\Phi(\gamma)$，这就证明了 $\Phi$ 是从 $\Gamma$ 到
$\pi^R_1((X,\sigma);x)$ 的群同态。

为断定 $\Phi$ 是群同构，我们注意到：把 $f\in\pi^R_1((X,\sigma);x)$ 送到道路
$\gamma:=p\circ\zeta$ 的同伦类的映射（其中 $\zeta$ 是 $\tilde X$ 中从 $\tilde x$ 到
$f(\tilde x)$ 的任一道路）是 $\Phi$ 的逆映射（这里再次用到 $\tilde X$ 单连通，以说明
该映射良定义）。
\end{proof}

把像 $\alpha:\pi^R_1((X,\sigma);x)\to\Z/2\Z$ 这样的满群同态称为 $\Z/2\Z$ 的一个
\emph{增广}（augmentation）会带来方便。我们将写作
\[
\alpha:\ \pi^R_1((X,\sigma);x)\longrightarrow\Z/2\Z\longrightarrow1,
\]
以记住 $\alpha$ 是满群同态。两个增广 $\alpha:A\to\Z/2\Z\to1$ 与
$\beta:B\to\Z/2\Z\to1$ 之间的同态，是使图
\[
\begin{tikzcd}[column sep=large, row sep=large]
A \arrow[r, "\alpha"] \arrow[d, "\psi"'] & \Z/2\Z \arrow[r] \arrow[d, "="] & 1\\
B \arrow[r, "\beta"'] & \Z/2\Z \arrow[r] & 1
\end{tikzcd}
\]
可交换的群同态 $\psi:A\to B$。特别地，$\psi$ 把 $\ker\alpha$ 送入 $\ker\beta$。用这套
语言说：实曲线 $(X,\sigma)$ 的基本群是 $\R$ 的绝对 Galois 群的增广。于是，两条实曲线
$(X',\sigma')$ 与 $(X,\sigma)$ 的基本群 $\pi^R_1((X',\sigma');x')$ 与
$\pi^R_1((X,\sigma);x)$ 之间的同态，就是使图
\[
\begin{tikzcd}[column sep=large, row sep=large]
1 \arrow[r] & \pi_1(X';x') \arrow[r] \arrow[d, "\psi\vert_{\pi_1(X';x')}"'] & \pi^R_1((X',\sigma');x') \arrow[r] \arrow[d, "\psi"'] & \Z/2\Z \arrow[r] \arrow[d, "="] & 1\\
1 \arrow[r] & \pi_1(X;x) \arrow[r] & \pi^R_1((X,\sigma);x) \arrow[r] & \Z/2\Z \arrow[r] & 1
\end{tikzcd}
\]
可交换的群同态 $\psi:\pi^R_1((X',\sigma');x')\to\pi^R_1((X,\sigma);x)$。利用推论 2.2.14
与命题 2.2.17 给出的实曲线基本群的描述，$\pi^R_1$ 的函子性相当容易建立。

\begin{proposition}\label{prop:2.2.18}
设 $(X',\sigma')$ 与 $(X,\sigma)$ 是两条实曲线，$\phi:X'\to X$ 是任一满足
$\phi\circ\sigma'=\sigma\circ\phi$ 的连续映射。取 $x'\in X'$，令
$x=\phi(x')\in X$。则存在增广的同态
\[
\begin{tikzcd}[column sep=large, row sep=large]
\pi^R_1((X',\sigma');x') \arrow[r] \arrow[d, "\phi_*"'] & \Z/2\Z \arrow[r] \arrow[d, "="] & 1\\
\pi^R_1((X,\sigma);x) \arrow[r] & \Z/2\Z \arrow[r] & 1
\end{tikzcd}
\]
特别地，实曲线的基本群是拓扑不变量。
\end{proposition}

\begin{proof}
由于条件 $\phi\circ\sigma'=\sigma\circ\phi$，若 $x'\in(X')^{\sigma'}$ 则 $x\in X^\sigma$，
故 $\pi^R_1((X',\sigma');x')$ 与 $\pi^R_1((X,\sigma);x)$ 都同构于推论 2.2.14 中那种形式
的半直积，而 $\phi_*$ 可以定义为把 $(\gamma,\varepsilon)\in
\pi_1(X';x')\rtimes_{\sigma'}\Z/2\Z$ 送到 $(\phi(\gamma),\varepsilon)\in
\pi_1(X;x)\rtimes_\sigma\Z/2\Z$。若现在 $x'\notin(X')^{\sigma'}$，则用命题 2.2.17
如下定义 $\phi_*$：若 $\sigma(x)\ne x$，把 $\gamma$（从 $x'$ 到 $x'$ 或从 $x'$ 到
$\sigma'(x')$）送到 $\phi(\gamma)$；否则，若 $\sigma(x)=x$，则当 $\gamma(1)=x'$ 时把
$\gamma$ 送到 $(\phi(\gamma),\Id_X)\in\pi_1(X;x)\rtimes_\sigma\Z/2\Z$，当
$\gamma(1)=\sigma'(x')$ 时送到 $(\phi(\gamma),\sigma)$。在所有情形下，$\phi_*$ 都是
$\Z/2\Z$ 的增广的同态。
\end{proof}

不借助道路来定义 $\phi_*$ 要稍难一些，但我们仍可如下勾勒函子性的证明。给定连续实同态
$\phi:(X',\sigma')\to(X,\sigma)$ 与实覆盖 $q:(Y,\tau)\to(X,\sigma)$，存在实覆盖
$\phi^*q:(\phi^*Y,\phi^*\tau)\to(X',\sigma')$，满足
\[
(\phi^*Y)_{x'}=Y_{\phi(x')}\xrightarrow{\ \phi^*\tau\ }(\phi^*Y)_{\sigma'(x')}
=Y_{\sigma(\phi(x'))}.
\]
特别地，集 $(\phi^*Y)_{x'}$ 的任一自同构都定义了集 $Y_{\phi(x')}=Y_x$ 的自同构，且与
实覆盖的同态相容。由于 $\pi^R_1((X',\sigma');x')$ 是 $(X',\sigma')$ 的一切实覆盖在
$x'$ 处纤维的自同构群，而 $\pi^R_1((X,\sigma);x)$ 是 $(X,\sigma)$ 的一切实覆盖在
$x=\phi(x')$ 处纤维的自同构群，上面的构造定义了映射
$\phi_*:\pi^R_1((X',\sigma');x')\to\pi^R_1((X,\sigma);x)$，它是 $\Z/2\Z$-增广的同态。

紧连通克莱因曲面基本群的展示式可参看 [9]；我们用两个简单的例子结束本节。

\begin{example}\label{ex:2.2.19}
若 $(X,\sigma)$ 是无实点的克莱因曲面，则
\[
\pi^R_1((X,\sigma);x)\simeq\pi_1^{\mathrm{top}}(X/\sigma;[x]).
\]
事实上，此时典范投影 $q:X\to X/\sigma$ 是覆盖映射，而 $\tilde X$ 是 $X/\sigma$ 的万有
覆盖空间。此外，自映射 $f:\tilde X\to\tilde X$ 覆盖 $X/\sigma$ 上的 $\Id_{X/\sigma}$，
当且仅当它覆盖 $X$ 上的 $\Id_X$ 或 $\sigma$，而这恰好意味着
$\pi_1^{\mathrm{top}}(X/\sigma;[x])\simeq\pi^R_1((X,\sigma);x)$。看出这一点的另一种
说法是：由拓扑覆盖 $q$ 定义的拉回映射 $q^*$ 是 $X/\sigma$ 的拓扑覆盖范畴与 $(X,\sigma)$
的实覆盖范畴之间的范畴等价，故两个 Galois 群同构。
\end{example}

\begin{example}\label{ex:2.2.20}
考虑非紧克莱因曲面
$X:=\CP^1\setminus\{s_1,s_2,s_3\}$，其中 $s_1,s_2,s_3$ 是
$\RP^1=\Fix(\sigma)\subset\CP^1$ 中互不相同的点；取 $x\notin\RP^1$ 为基点。用 $\eta_1$
记从 $x$ 到 $\sigma(x)$、穿过 $s_1$ 与 $s_2$ 之间且不绕任何穿刺点（puncture）的道路的同伦类；
同样，$\eta_2$ 穿过 $s_2$ 与 $s_3$ 之间，$\eta_3$ 穿过 $s_3$ 与 $s_1$ 之间。于是每个
$\eta_i$ 在 $\pi^R_1((X,\sigma);x)$ 中都是二阶元，而
$\zeta_1:=\eta_2\boxtimes\eta_1$、$\zeta_2:=\eta_3\boxtimes\eta_2$、
$\zeta_3:=\eta_1\boxtimes\eta_3$ 属于 $\pi_1(X;x)$ 且满足
$\zeta_3*\zeta_2*\zeta_1=1$。容易验证 $\zeta_1,\zeta_2,\zeta_3$ 生成 $\pi_1(X;x)$，
$\eta_1,\eta_2,\eta_3$ 生成 $\pi^R_1((X,\sigma);x)$。更精确地说，我们有下述两个展示：
\[
\pi_1(X;x)=\langle a_1,a_2,a_3\mid a_1a_2a_3=1\rangle,
\qquad
\pi^R_1((X,\sigma);x)=\langle b_1,b_2,b_3\mid b_1^2=b_2^2=b_3^2=1\rangle.
\]
增广映射 $\pi^R_1((X,\sigma);x)\to\Z/2\Z$ 把每个 $b_i$ 送到 $\Z/2\Z$ 的非平凡元，
故元素 $a_1=b_1b_2$、$a_2=b_2b_3$、$a_3=b_3b_1$ 都被送到平凡元。序列
\[
1\longrightarrow\langle b_1b_2,\,b_2b_3,\,b_3b_1\rangle
\longrightarrow\langle b_1,b_2,b_3\mid b_1^2=b_2^2=b_3^2=1\rangle
\longrightarrow\Z/2\Z\longrightarrow1
\]
就是此情形下 $\pi^R_1((X,\sigma);x)$ 的同伦正合列。本例以直接的方式推广到
$s_1,\dots,s_n\in\RP^1$ 中任意有限多个穿刺点。注意：若我们不要求 $s_i$ 落在
$\RP^1$ 中，则由于 $X$ 上实结构的存在，不在 $\RP^1$ 中的那些必须成对出现
$(s_i,\sigma(s_i))$。
\end{example}
